% !Mode:: "TeX:UTF-8"

\hitsetup{
  %******************************
  % 注意：
  %   1. 配置里面不要出现空行
  %   2. 不需要的配置信息可以删除
  %******************************
  %
  %=====
  % 秘级
  %=====
  statesecrets={公开},
  natclassifiedindex={TM301.2},
  intclassifiedindex={62-5},
  %
  %=========
  % 中文信息
  %=========
  ctitlecover={面向多泛化场景的模型迁移研究},%放在封面中使用，自由断行
  ctitle={面向多泛化场景的模型迁移研究},%放在原创性声明中使用
  % csubtitle={一条副标题}, %一般情况没有，可以注释掉
  cxueke={工学},
  csubject={计算机科学与技术},
  caffil={计算学部},
  cauthor={陈希},
  csupervisor={姚鸿勋教授},
  % cassosupervisor={某某某教授}, % 副指导老师
  % ccosupervisor={某某某教授}, % 联合指导老师
  % 如果是第一封面的日期要手动设置，需要取消注释下一行，并将内容改为“规范”中要求的封面第一页最下方的日期
  % firstpagecdate={2022年6月},
  % 日期自动使用当前时间，若需指定按如下方式修改：
  % cdate={超新星纪元},
  cstudentid={9527},
  cstudenttype={学术学位论文}, %非全日制教育申请学位者
  cnumber={no9527}, %编号
  cpositionname={哈铁西站}, %博士后站名称
  cfinishdate={20XX年X月---20XX年X月}, %到站日期
  csubmitdate={20XX年X月}, %出站日期
  cstartdate={3050年9月10日}, %到站日期
  cenddate={3090年10月10日}, %到站日期
  %（同等学力人员）、（工程硕士）、（工商管理硕士）、
  %（高级管理人员工商管理硕士）、（公共管理硕士）、（中职教师）、（高校教师）等
  %
  %
  %=========
  % 英文信息
  %=========
  etitle={Research on Model Transfer in Multi-Generalization Scenarios},
  esubtitle={This is the sub title},
  exueke={Engineering},
  esubject={Computer Science and Technology},
  eaffil={\emultiline[t]{Faculty of Computing}},
  eauthor={Xi Chen},
  esupervisor={Prof. Hongxun Yao},
  % eassosupervisor={XXX},
  % 日期自动生成，若需指定按如下方式修改：
  % edate={December, 2017},
  estudenttype={Master of Art},
  %
  % 关键词用“英文逗号”分割
  ckeywords={模型迁移；多泛化场景；样本挖掘；多模态监督；语义理解标准},
  ekeywords={Model Transfer; Multi-Generalization Scenarios; Sample Mining; Multimodal Supervision; Semantic Understanding Standards},
}

\begin{cabstract}

视觉感知模型从封闭实验环境走向真实开放场景的过程中，需要应对训练与部署之间多层面的差异。视觉感知任务在预训练—微调范式下已取得显著进展，然而大规模预训练模型在训练阶段所接触的样本分布、类别词表与语义定义，往往与目标部署环境存在系统性偏差：这种偏差不仅体现在数据的统计特性上，也体现在类别集合的开放程度与语义刻画的精细程度上。当模型从实验室走向真实环境时，缺乏目标标注条件下的分布差异持续存在、已知与未见类别在开放环境中共存、高阶语义因主观多义而难以仅由结论判定等问题相互交织，使得模型的可靠迁移成为制约感知系统实际落地的核心瓶颈。本文将此类由训练与部署条件不一致引发的稳健性问题概括为多泛化场景，并依据偏移发生的空间层次将其划分为域泛化、类泛化与语义泛化三个层面，三者依次对应特征与样本空间、类别空间与语义空间的偏移，围绕不同偏移维度、监督资源约束和模型载体，研究模型知识向目标任务的可靠迁移与可诊断评测。

围绕多泛化场景下的模型迁移，本文沿样本表征、类别概念与高阶语义三个层次展开研究。面向数据分布变化的域泛化场景，研究无监督与少监督条件下的目标域适应；面向类别词表开放的类泛化场景，研究未见类别的识别与像素级定位；面向语义理解要求提升的任务场景，构建分层评测协议，考察多模态大语言模型能否依据视觉证据作出情感解释。前两类研究着重提升模型的迁移能力，后一类研究刻画已有能力在高阶语义任务中的适用范围。主要研究内容如下：

首先，面向跨域分布偏移下的无监督模型迁移场景，模型既要在目标域维持稳定的判别表征，也要避免错误预测在自训练中被反复强化。现有伪标签方法常依据预测置信度或固定阈值选择样本，但域偏移下的高置信预测仍可能出错；与此同时，骨干网络对源域纹理、背景及空间共现的依赖可能削弱目标表征的稳定性。在源数据不可访问、目标标签缺失的设定下，这两类不确定性更难借助真实标签或源样本加以校正。为此，本文沿监督信号与表征建模两条路径展开研究：提出基于特征空间邻居一致性的自适应伪标签筛选方法，以邻居分布估计样本级可靠性，并通过类别感知对比学习抑制噪声传播；进一步提出基于选择性扫描的不确定性感知迁移方法，利用通道级视觉状态空间模块与语义一致性打乱策略，缓解通道建模不足和虚假相关性累积，以线性复杂度实现长程跨域建模。两项研究在 Office-Home、VisDA-C、Office、DomainNet-126 四个基准上分别将精度提升至 81.7\%、89.3\%、93.3\%、77.9\%，取得了具有竞争力的迁移性能。

其次，面向少监督模型迁移场景，有限监督的核心挑战在于如何提高每一份目标域标注的学习收益，使模型在分布偏移下仍获得可靠判别能力。现有方法分别着力于有限标注的获取或利用，尚难同时兼顾目标样本的覆盖性、不同知识来源的互补性，以及监督信息向目标模型能力的转化。视觉—语言模型能够提供与源模型互补的开放语义先验，但其潜力仍取决于目标域适配与有效的跨模型交互。在主动标注预算有限的具体设定下，本文提出基于多模态基础模型引导的双主动学习方法，通过双聚焦样本选择、视觉—语言模型适配与交替蒸馏，促进源模型与视觉—语言模型之间的双向知识传递，并以协作风险上界分析样本覆盖、模型风险及跨模型对齐对目标性能的影响。在 Office-Home、VisDA-C、Office 闭集设定下，该方法较最优基线分别提升 4.5、1.5、0.6 个百分点，在开集设定下分别提升 9.7、3.5、1.0 个百分点。

再次，面向开放词汇分割任务，模型既要识别训练阶段未见的类别，也要准确界定其像素区域。现有两阶段方法通常先生成类别无关掩码、再进行语义分类，掩码与类别判别所依赖的特征难以充分对齐；单阶段方法虽简化流程，仍受图文对齐模型缺乏像素级监督的限制。分割基础模型具有较强的空间定位能力，却不能独立提供开放词汇语义。针对两类先验在掩码质量与类别判别上的能力错配，本文提出跨模型先验融合方法，在轻量解码器中通过查询与特征注入融合空间定位和类别语义信息，并在推理阶段集成分割基础模型的零样本掩码预测，以改善未见类别的分割质量。该方法在全景分割、语义分割和实例分割等八个基准上验证了其泛化能力。

最后，面向艺术情感推理任务，目标任务从事实识别延伸到情感归因与证据化解释，模型生成的流畅描述并不必然具有可靠的视觉依据。现有视觉情感数据集和美学评测多关注情感标签、属性评分或描述质量，难以区分模型是否识别了支撑情感判断的感知线索，也难以定位从感知到推理的具体失效环节。针对这一评测缺口，本文构建多层级基准 AER-Bench，将艺术情感理解拆解为感知维度、情感与效价—唤醒—支配、基于证据的推理、情感描述四个互补层级，并结合结构化题目与开放式评分进行诊断。该基准覆盖 2{,}058 件艺术作品、34{,}791 道评测题目；对 14 个代表性多模态大语言模型的系统评测揭示了情感判断与证据化解释之间的能力差距。

综上，本文围绕多泛化场景下的模型迁移，形成了从跨域稳定表征、开放类别识别到高阶语义能力评测的研究路径：通过无监督与少监督适应提升跨分布迁移的可靠性，通过跨模型先验融合增强未见类别的细粒度识别，并以多维评测协议刻画多模态大语言模型在证据化语义理解中的能力边界。

\end{cabstract}
\begin{eabstract}

  As visual perception models move from closed experimental environments to open real-world scenarios, they must cope with multifaceted discrepancies between training and deployment. Visual perception tasks have achieved remarkable progress under the pre-training--fine-tuning paradigm, yet the sample distributions, category vocabularies, and semantic definitions encountered by large-scale pretrained models during training often differ systematically from those in target deployment environments: such discrepancies are manifested not only in the statistical properties of data, but also in the openness of category sets and the granularity of semantic characterization. When models are transferred from laboratories to real-world settings, distributional discrepancies that persist in the absence of target annotations, the coexistence of known and unseen categories in open environments, and high-level semantics that cannot be judged by the final prediction alone become intertwined, making reliable model transfer a central bottleneck for the practical deployment of perception systems. This dissertation refers to such robustness issues arising from inconsistencies between training and deployment conditions as \emph{multi-generalization scenarios}, and partitions them according to the spatial level at which the shift occurs into three interlocking levels of domain, category, and semantic generalization, which correspond in turn to shifts in feature and sample space, category space, and semantic space. The reliable transfer and diagnosable evaluation of model knowledge for target tasks are investigated under different shift dimensions, supervision-resource constraints, and model carriers.
  
  Taking model transfer under multi-generalization scenarios as its research line, this dissertation focuses on three generalization challenges arising in the transition from training to deployment, namely domain distribution shift, the coexistence of open-set categories, and high-level semantic understanding, and organizes the research of the dissertation accordingly. At the domain-generalization level, this dissertation studies how unlabeled target data can be exploited reliably under constrained supervision resources, and develops both unsupervised and weakly supervised transfer methods. At the category-generalization level, this dissertation studies the semantic coverage of models when the category set is no longer closed, and develops open-vocabulary segmentation methods for unseen categories. At the semantic-generalization level, this dissertation studies whether the model's judgments of emotion and other semantics genuinely rely on the visual evidence in the image, and develops evaluation protocols with stratified diagnostic capability. The three levels progress in turn and jointly point toward the reliable deployment of visual perception models. The main research contents are as follows:
  
  First, to suppress uncertainty in unsupervised model transfer, this dissertation investigates the reliable quantification and control of cross-domain transfer uncertainty. To mitigate the accumulation of transfer uncertainty caused by pseudo-label noise, an adaptive pseudo-label filtering method based on neighborhood consistency is proposed. The method estimates sample-level uncertainty from neighborhood distributions to adaptively select high-quality pseudo-labels. To address two inherent limitations of state-space models, namely insufficient channel modeling and the accumulation of spurious correlations, an uncertainty-aware transfer method based on selective scanning is further proposed. By incorporating a channel-wise visual state-space module and a semantic-consistent shuffling strategy, the method enables efficient long-range cross-domain modeling with linear complexity. On the Office-Home, VisDA-C, Office, and DomainNet-126 benchmarks, the proposed methods achieve accuracies of 81.7\%, 89.3\%, 93.3\%, and 77.9\%, respectively, demonstrating competitive transfer performance.
  
  Second, this dissertation addresses information maximization in weakly supervised model transfer, namely maximizing the amount of information that the model acquires from the supervisory signals under a limited annotation budget. Existing methods lack systematic exploration of supervisory signals from multimodal pretrained models in both sample selection and model learning, fail to convert the complementary information across channels into learning gains of the target model, and have not established a rigorous theoretical framework for characterizing information content and for bidirectional deep collaboration between vision--language models and source models. To this end, a dual active learning method guided by multimodal foundation models is proposed. Through sample-selection design, multimodal model adaptation, and alternating distillation, the method achieves bidirectional deep collaboration between source models and vision--language models in the hypothesis space. On the closed-set settings of Office-Home, VisDA-C, and Office, it improves upon the strongest baselines by 4.5, 1.5, and 0.6 percentage points, respectively; on the open-set settings, the corresponding improvements are 9.7, 3.5, and 1.0 percentage points.
  
  Third, to deeply exploit multimodal category priors for model generalization, this dissertation investigates multi-stream semantic alignment and spatial-localization reasoning optimization. Image--text alignment models tend to emphasize global semantics while producing coarse mask boundaries, whereas foundation segmentation models excel at localization but lack semantic generalization. Moreover, these two types of priors are difficult to deeply integrate within a unified decoding framework. To address this capability mismatch, a cross-model prior fusion method is proposed. It achieves deep fusion of spatial-localization and semantic features within a lightweight decoder, and incorporates zero-shot mask predictions from a foundation segmentation model through an open-set ensemble module during inference, thereby improving the segmentation quality of unseen categories. The proposed method validates its generalization capability for unseen categories on eight benchmarks spanning panoptic, semantic, and instance segmentation.
  
  Finally, to establish multidimensional standards for semantic understanding, this dissertation studies a multi-level evaluation protocol for art emotion reasoning, enabling distinguishable diagnosis between fluent emotional descriptions and their reliable grounding in visual evidence in higher-level semantic scenarios. To address the lack of fine-grained diagnostic capability in existing visual-emotion benchmarks for differentiating fluent emotional expression from evidence-grounded explanation, AER-Bench, a multi-level benchmark for art emotion reasoning, is constructed. It decomposes art emotion understanding into four complementary levels: perceptual dimensions; emotion together with valence, arousal, and dominance; evidence-grounded reasoning; and emotion description. AER-Bench contains 2{,}058 artworks and 34{,}791 evaluation questions, and has been used to systematically evaluate 14 representative multimodal large language models.
  
  In summary, this dissertation systematically investigates model transfer in multi-generalization scenarios from the perspectives of domain, category, and semantic generalization. At the domain-generalization level, it studies unsupervised and weakly supervised model transfer: the former mitigates transfer uncertainty by suppressing noisy pseudo-labels and their propagation, whereas the latter leverages open-semantic priors from multimodal models to drive sample mining and cross-modal collaborative transfer. At the category-generalization level, it investigates multi-stream semantic alignment and spatial-localization reasoning optimization, integrating the localization and semantic priors of heterogeneous models to improve fine-grained recognition of unseen categories. At the semantic-generalization level, it establishes multidimensional semantic-understanding standards by constructing evaluation protocols that jointly support evidence anchoring and multidimensional diagnosis, thereby providing interpretable evaluation criteria for open semantic scenarios.
  
  \end{eabstract}
